\chapter{2011年重庆卷（文）}

\section{单选题}

\begin{problem}
在等差数列 \(\{a_{n}\}\) 中，\(a_{2} = 2\)，\(a_{3} = 4\)，则 \(a_{10} =\)（\quad）

\choices
    {\(12\)}
    {\(14\)}
    {\(16\)}
    {\(18\)}
\end{problem}

\begin{answer}
    D
\end{answer}

\begin{solution}
等差数列的公差为 \(d=a_{3}-a_{2}=4-2=2\)，所以
\(a_{10}=a_{2}+8d=2+8\cdot2=18\). 故选 D.
\end{solution}

\begin{problem}
设 \(U = \R\)，\(M = \{x \mid x^2 - 2x > 0\}\)，则 \(\complement_U M =\)（\quad）

\choices
    {\(\left[0,2\right]\)}
    {\(\left(0,2\right)\)}
    {\(\left(-\infty,0\right)\cup\left(2,+\infty\right)\)}
    {\(\left(-\infty,0\right]\cup\left[2,+\infty\right)\)}
\end{problem}

\begin{answer}
    A
\end{answer}

\begin{solution}
由 \(x^{2}-2x=x(x-2)>0\)，得 \(x<0\) 或 \(x>2\)，所以
\(M=(-\infty,0)\cup(2,+\infty)\). 取补集可得
\(\complement_{U}M=[0,2]\). 故选 A.
\end{solution}

\begin{problem}
曲线 \( y = -x^{3} + 3x^{2} \) 在点\(\left(1,2\right)\)处的切线方程为（\quad）

\choices
    {\(y = 3x - 1\)}
    {\(y = -3x + 5\)}
    {\(y = 3x + 5\)}
    {\(y = 2x\)}
\end{problem}

\begin{answer}
    A
\end{answer}

\begin{solution}
令 \(y=-x^{3}+3x^{2}\)，则 \(y'=-3x^{2}+6x\)，从而曲线在 \(x=1\) 处的切线斜率为 \(y'(1)=3\). 切线过点 \((1,2)\)，故
\(y-2=3(x-1)\)，即 \(y=3x-1\). 故选 A.
\end{solution}

\begin{problem}
从一堆苹果中任取 \(10\) 只，称得它们的质量如下（单位：\(\text{克}\)）：\(125\)，\(120\)，\(122\)，\(105\)，\(130\)，\(114\)，\(116\)，\(95\)，\(120\)，\(134\)，则样本数据落在 \(\left[114.5,124.5\right]\) 内的频率为（\quad）

\choices
    {\(0.2\)}
    {\(0.3\)}
    {\(0.4\)}
    {\(0.5\)}
\end{problem}

\begin{answer}
    C
\end{answer}

\begin{solution}
区间 \([114.5,124.5]\) 中的数据为 \(120,122,116,120\)，共有 \(4\) 个. 样本总数为 \(10\)，故频率为
\(\frac{4}{10}=0.4\). 故选 C.
\end{solution}

\begin{problem}
已知向量 \(\bs{a} = (1, k)\)，\(\bs{b} = (2, 2)\)，且 \(\bs{a} + \bs{b}\) 与 \(\bs{a}\) 共线，那么 \(\bs{a} \cdot \bs{b}\) 的值为（\quad）

\choices
    {\(1\)}
    {\(2\)}
    {\(3\)}
    {\(4\)}
\end{problem}

\begin{answer}
    D
\end{answer}

\begin{solution}
\(\bs{a}+\bs{b}=(3,k+2)\) 与 \(\bs{a}=(1,k)\) 共线，所以
\(1\cdot(k+2)-3k=0\)，解得 \(k=1\). 因此
\(\bs{a}\cdot\bs{b}=1\cdot2+1\cdot2=4\). 故选 D.
\end{solution}

\begin{problem}
设 \(a = \log_{\frac{1}{3}} \frac{1}{2}\)，\(b = \log_{\frac{1}{3}} \frac{2}{3}\)，\(c = \log_3 \frac{4}{3}\)，则 \(a\)，\(b\)，\(c\) 的大小关系是（\quad）

\choices
    {\(a < b < c\)}
    {\(c < b < a\)}
    {\(b < a < c\)}
    {\(b < c < a\)}
\end{problem}

\begin{answer}
    B
\end{answer}

\begin{solution}
令 \(r=\log_{3}2\)，则
\(a=\log_{1/3}(1/2)=r\)，
\(b=\log_{1/3}(2/3)=1-r\)，
\(c=\log_{3}(4/3)=2r-1\). 由于 \(2>\sqrt{3}\)，所以 \(r>\frac12\)，从而 \(b<a\)；又由于 \(2^{3}<3^{2}\)，所以 \(r<\frac23\)，从而 \(c<b\). 因此 \(c<b<a\). 故选 B.
\end{solution}

\begin{problem}
若函数 \( f(x) = x + \frac{1}{x - 2} \)（\(x > 2\)）在 \( x = a \) 处取最小值，则 \( a = \)（\quad）

\choices
    {\(1 + \sqrt{2}\)}
    {\(1 + \sqrt{3}\)}
    {\(3\)}
    {\(4\)}
\end{problem}

\begin{answer}
    C
\end{answer}

\begin{solution}
令 \(t=x-2>0\)，则
\(f(x)=x+\frac{1}{x-2}=t+2+\frac1t\). 由基本不等式
\(t+\frac1t\geq2\)，等号当且仅当 \(t=1\) 时成立，所以 \(f(x)\geq4\)，且最小值在 \(x=t+2=3\) 处取得. 故选 C.
\end{solution}

\begin{problem}
若 \(\triangle ABC\) 的内角 \(A\)，\(B\)，\(C\) 满足 \(6\sin A = 4\sin B = 3\sin C\)，则 \(\cos B =\)（\quad）

\choices
    {\(\frac{\sqrt{15}}{4}\)}
    {\(\frac{3}{4}\)}
    {\(\frac{3\sqrt{15}}{16}\)}
    {\(\frac{11}{16}\)}
\end{problem}

\begin{answer}
    D
\end{answer}

\begin{solution}
设三角形 \(ABC\) 的边 \(a\)，\(b\)，\(c\) 分别对着角 \(A\)，\(B\)，\(C\). 由正弦定理及题设，
\(a:b:c=\sin A:\sin B:\sin C=2:3:4\)，故可设 \(a=2k\)，\(b=3k\)，\(c=4k\)，其中 \(k>0\). 由余弦定理，
\[
\cos B=\frac{a^{2}+c^{2}-b^{2}}{2ac}=\frac{(2k)^{2}+(4k)^{2}-(3k)^{2}}{2\cdot2k\cdot4k}=\frac{11}{16}
\]
故选 D.
\end{solution}

\begin{problem}
设双曲线的左准线与两条渐近线交于 \(A\)，\(B\) 两点，左焦点在以 \(AB\) 为直径的圆内，则该双曲线的离心率取值范围为（\quad）

\choices
    {\(\left(0,\sqrt{2}\right)\)}
    {\(\left(1,\sqrt{2}\right)\)}
    {\(\left(\frac{\sqrt{2}}{2}, 1\right)\)}
    {\(\left(\sqrt{2},+\infty\right)\)}
\end{problem}

\begin{answer}
    B
\end{answer}

\begin{solution}
按横轴为实轴设双曲线为 \(\frac{x^{2}}{a^{2}}-\frac{y^{2}}{b^{2}}=1\)，其中 \(a\)，\(b>0\)，并记 \(c=\sqrt{a^{2}+b^{2}}\)，离心率 \(e=\frac ca\). 左准线为 \(x=-\frac{a^{2}}c\)，两条渐近线为 \(y=\pm\frac ba x\)，所以它们的交点为
\(A=(-\frac{a^{2}}c,\frac{ab}c)\) 与 \(B=(-\frac{a^{2}}c,-\frac{ab}c)\). 以 \(AB\) 为直径的圆的圆心为 \(E=(-\frac{a^{2}}c,0)\)，半径为 \(\frac{ab}c\). 左焦点为 \(F=(-c,0)\)，故
\(EF=c-\frac{a^{2}}c=\frac{b^{2}}c\).

左焦点在圆内等价于 \(EF<\frac{ab}c\)，即 \(b<a\). 于是
\(c=\sqrt{a^{2}+b^{2}}\) 满足 \(a<c<\sqrt2a\)，从而 \(1<e<\sqrt2\). 故选 B.
\end{solution}

\begin{problem}
高为 \(\sqrt{2}\) 的四棱锥 \(S - ABCD\) 的底面是边长为 \(1\) 的正方形，点 \(S\)，\(A\)，\(B\)，\(C\)，\(D\) 均在半径为 \(1\) 的同一球面上，则底面 \(ABCD\) 的中心与顶点 \(S\) 之间的距离为（\quad）

\choices
    {\(\frac{\sqrt{10}}{2}\)}
    {\(\frac{\sqrt{2} + \sqrt{3}}{2}\)}
    {\(\frac{3}{2}\)}
    {\(\sqrt{2}\)}
\end{problem}

\begin{answer}
    A
\end{answer}

\begin{solution}
设正方形底面的中心为 \(O\)，底面所在平面为 \(z=0\)，球心为 \(E\). 正方形边长为 \(1\)，所以 \(OA=\frac{\sqrt2}{2}\). 球心到四个顶点等距，故 \(E\) 在过 \(O\) 且垂直于底面的直线上. 由球半径为 \(1\)，
\(OE^{2}+OA^{2}=1\)，得 \(OE=\frac{\sqrt2}{2}\).

设 \(H\) 为 \(S\) 在底面上的射影. 若 \(E\) 在底面下方，则 \(S\) 与 \(E\) 的竖直距离为 \(\sqrt2+\frac{\sqrt2}{2}>1\)，不可能同在半径为 \(1\) 的球面上，因此 \(E\) 在底面上方. 于是 \(S\) 与 \(E\) 的竖直距离为 \(\sqrt2-\frac{\sqrt2}{2}=\frac{\sqrt2}{2}\)，由 \(SE=1\) 得
\(OH^{2}=1-\left(\frac{\sqrt2}{2}\right)^{2}=\frac12\). 因此
\(OS^{2}=OH^{2}+SH^{2}=\frac12+2=\frac52\)，即 \(OS=\frac{\sqrt{10}}2\). 故选 A.
\end{solution}

\section{填空题}

\begin{problem}
\(\left(1 + 2x\right)^{6}\) 的展开式中 \(x^{4}\) 的系数是\fillinblank{}.
\end{problem}

\begin{answer}
    \(240\)
\end{answer}

\begin{solution}
由二项式定理，\(x^{4}\) 项取四个 \(2x\) 及两个 \(1\)，所以其系数为
\(\mathrm C_{6}^{4}\cdot2^{4}=15\cdot16=240\).
\end{solution}

\begin{problem}
若 \(\cos \alpha = -\frac{3}{5}\)，且 \(\alpha \in \left(\pi, \frac{3\pi}{2}\right)\)，则 \(\tan \alpha =\)\fillinblank{}.
\end{problem}

\begin{answer}
    \(\frac{4}{3}\)
\end{answer}

\begin{solution}
因为 \(\alpha\in(\pi,\frac{3\pi}{2})\)，所以 \(\sin\alpha<0\). 由
\(\sin^{2}\alpha=1-\cos^{2}\alpha=1-\frac{9}{25}=\frac{16}{25}\)，得 \(\sin\alpha=-\frac45\). 因此
\(\tan\alpha=\frac{\sin\alpha}{\cos\alpha}=\frac{-4/5}{-3/5}=\frac43\).
\end{solution}

\begin{problem}
过原点的直线与圆 \(x^{2} + y^{2} - 2x - 4y + 4 = 0\) 相交所得弦的长为\(2\)，则该直线的方程为\fillinblank{}.
\end{problem}

\begin{answer}
    \(2x-y=0\)
\end{answer}

\begin{solution}
将圆的方程配方得
\((x-1)^{2}+(y-2)^{2}=1\)，所以圆心为 \((1,2)\)，半径为 \(1\). 弦长为 \(2\) 时，该弦必须是直径，因此过原点的直线必须经过圆心 \((1,2)\). 过原点和 \((1,2)\) 的直线为 \(y=2x\)，即
\(2x-y=0\).
\end{solution}

\begin{problem}
从甲，乙等\(10\)位同学中任选\(3\)位去参加某项活动，则所选\(3\)位中有甲但没有乙的概率为\fillinblank{}.
\end{problem}

\begin{answer}
    \(\frac{7}{30}\)
\end{answer}

\begin{solution}
从 \(10\) 位同学中任选 \(3\) 位的总数为 \(\mathrm C_{10}^{3}\). 有甲但没有乙时，固定选甲，不选乙，再从其余 \(8\) 位中选 \(2\) 位，符合条件的选法数为 \(\mathrm C_{8}^{2}\). 因此所求概率为
\(\frac{\mathrm C_{8}^{2}}{\mathrm C_{10}^{3}}=\frac{28}{120}=\frac{7}{30}\).
\end{solution}

\begin{problem}
若实数 \(a\)，\(b\)，\(c\) 满足 \(2^{a} + 2^{b} = 2^{a + b}\)，\(2^{a} + 2^{b} + 2^{c} = 2^{a + b + c}\)，则 \(c\) 的最大值是\fillinblank{}.
\end{problem}

\begin{answer}
    \(2-\log_{2}3\)
\end{answer}

\begin{solution}
令 \(x=2^{a}\)，\(y=2^{b}\)，\(z=2^{c}\)，则 \(x\)，\(y\)，\(z>0\)，且题设化为
\(x+y=xy\)，\(x+y+z=xyz\). 第一式等价于 \((x-1)(y-1)=1>0\)，所以 \(x-1\) 与 \(y-1\) 同号. 因为 \(x\)，\(y>0\)，若 \(0<x<1\) 且 \(0<y<1\)，则 \(0<(x-1)(y-1)<1\)，不可能成立，所以 \(x\)，\(y>1\). 令 \(s=x+y=xy\)，则
\(s=2+(x-1)+(y-1)\geq4\)，等号当且仅当 \(x-1=y-1=1\)，即 \(x=y=2\).

由第二式得 \(s+z=sz\)，所以
\(z=\frac{s}{s-1}=1+\frac1{s-1}\leq\frac43\). 因此
\(c=\log_{2}z\leq\log_{2}\frac43=2-\log_{2}3\). 当 \(x=y=2\)、\(z=\frac43\) 时等号成立，故 \(c\) 的最大值为 \(2-\log_{2}3\).
\end{solution}

\section{解答题}

\begin{problem}
设 \(\{a_{n}\}\) 是公比为正数的等比数列，\(a_{1} = 2\)，\(a_{3} = a_{2} + 4\).
\begin{enumerate}
\item 求 \( \{a_{n}\} \) 的通项公式；
\item 设 \( \{b_{n}\} \) 是首项为 1，公差为 2 的等差数列. 求数列 \( \{a_{n} + b_{n}\} \) 的前 \(n\) 项和 \(S_{n}\).
\end{enumerate}
\end{problem}

\begin{solution}
\begin{enumerate}
\item 设等比数列的公比为 \(q>0\). 由 \(a_{1}=2\) 得 \(a_{2}=2q\)，\(a_{3}=2q^{2}\). 代入 \(a_{3}=a_{2}+4\)，得
\(2q^{2}=2q+4\)，即 \(q^{2}-q-2=0\). 所以 \(q=2\) 或 \(q=-1\)，由 \(q>0\) 舍去 \(q=-1\)，从而
\(a_{n}=a_{1}q^{n-1}=2\cdot2^{n-1}=2^{n}\).
\item 等差数列 \(\{b_{n}\}\) 的通项为
\(b_{n}=1+2(n-1)=2n-1\). 因此
\(a_{n}+b_{n}=2^{n}+2n-1\)，其前 \(n\) 项和为
\[
S_{n}=\sum_{k=1}^{n}(2^{k}+2k-1)
=\sum_{k=1}^{n}2^{k}+\sum_{k=1}^{n}(2k-1)
=(2^{n+1}-2)+(n(n+1)-n)
=2^{n+1}+n^{2}-2
\]
\end{enumerate}
\end{solution}

\begin{problem}
某市公租房的房源位于 \(A\)，\(B\)，\(C\) 三个片区. 设每位申请人只申请其中一个片区的房源，且申请其中任一个片区的房源是等可能的，求该市的任 \(4\) 位申请人中：
\begin{enumerate}
\item 没有人申请 \(A\) 片区房源的概率；
\item 每个片区的房源都有人申请的概率.
\end{enumerate}
\end{problem}

\begin{solution}
\begin{enumerate}
\item 每位申请人有 \(3\) 种等可能选择，\(4\) 位申请人的基本事件总数为 \(3^{4}\). 没有人申请 \(A\) 片区时，每人只能申请 \(B\) 或 \(C\)，共有 \(2^{4}\) 种，因此概率为
\(\frac{2^{4}}{3^{4}}=\frac{16}{81}\).
\item 每个片区都有人申请，用排除法计算. 总数为 \(3^{4}\). 指定一个片区无人申请时有 \(2^{4}\) 种，故至少有一个片区无人申请的计数为
\(\mathrm C_{3}^{1}2^{4}\). 若指定两个片区无人申请，则 \(4\) 位申请人只能申请剩下的一个片区，共 \(1^{4}\) 种；这类情况在前一计数中被重复计算，按容斥原理应加回 \(\mathrm C_{3}^{2}1^{4}\). 所以所求概率为
\[
\frac{3^{4}-\mathrm C_{3}^{1}2^{4}+\mathrm C_{3}^{2}1^{4}}{3^{4}}
=\frac{81-48+3}{81}
=\frac{36}{81}
=\frac49
\]
\end{enumerate}
\end{solution}

\begin{problem}
设函数 \(f(x) = \sin x \cos x - \sqrt{3}\cos\left(\pi + x\right)\cos x\) （\(x \in \R\)）.
\begin{enumerate}
\item 求 \(f(x)\) 的最小正周期；
\item 若函数 \(y = f(x)\) 的图象按 \(\bs{b} = \left(\frac{\pi}{4}, \frac{\sqrt{3}}{2}\right)\) 平移后得到函数 \(y = g(x)\) 的图象，求 \(y = g(x)\) 在 \(\left[0, \frac{\pi}{4}\right]\) 上的最大值.
\end{enumerate}
\end{problem}

\begin{solution}
\begin{enumerate}
\item 由 \(\cos(\pi+x)=-\cos x\)，得
\[
f(x)=\sin x\cos x+\sqrt3\cos^{2}x
=\frac12\sin2x+\frac{\sqrt3}{2}(1+\cos2x)
=\sin\left(2x+\frac{\pi}{3}\right)+\frac{\sqrt3}{2}
\]
因此 \(f(x)\) 的角频率为 \(2\)，其最小正周期为
\(T=\frac{2\pi}{2}=\pi\).
\item 图象按向量 \(\left(\frac{\pi}{4},\frac{\sqrt3}{2}\right)\) 平移后，
\[
g(x)=f\left(x-\frac{\pi}{4}\right)+\frac{\sqrt3}{2}
=\sin\left(2x-\frac{\pi}{6}\right)+\sqrt3
\]
当 \(x\in\left[0,\frac{\pi}{4}\right]\) 时，
\(2x-\frac{\pi}{6}\in\left[-\frac{\pi}{6},\frac{\pi}{3}\right]\). 该区间在 \([-\frac{\pi}{2},\frac{\pi}{2}]\) 内，正弦函数递增，所以最大值在 \(x=\frac{\pi}{4}\) 处取得，且
\[
g\left(\frac{\pi}{4}\right)=\sin\frac{\pi}{3}+\sqrt3=\frac{3\sqrt3}{2}
\]
\end{enumerate}
\end{solution}

\begin{problem}
设 \( f(x) = 2x^{3} + ax^{2} + bx + 1 \) 的导数为 \( f'(x) \). 若函数 \( y = f'(x) \) 的图象关于直线 \( x = -\frac{1}{2} \) 对称，且 \( f'(1) = 0 \).
\begin{enumerate}
\item 求实数 \(a\)，\(b\) 的值；
\item 求函数 \( f(x) \) 的极值.
\end{enumerate}
\end{problem}

\begin{solution}
\begin{enumerate}
\item
\(f'(x)=6x^{2}+2ax+b\) 的图象是开口向上的抛物线，其对称轴为
\(x=-\frac{2a}{2\cdot6}=-\frac{a}{6}\). 由题设对称轴为 \(x=-\frac12\)，得 \(a=3\). 再由 \(f'(1)=0\)，得
\(6+2\cdot3+b=0\)，所以 \(b=-12\).
\item 此时
\[
f'(x)=6x^{2}+6x-12=6(x+2)(x-1)
\]
所以 \(f'(x)>0\) 于 \((-\infty,-2)\) 和 \((1,+\infty)\)，\(f'(x)<0\) 于 \((-2,1)\). 因此 \(f\) 在 \(x=-2\) 处取得极大值，在 \(x=1\) 处取得极小值. 代入
\(f(x)=2x^{3}+3x^{2}-12x+1\)，得
\(f(-2)=-16+12+24+1=21\)，\(f(1)=2+3-12+1=-6\).
\end{enumerate}
\end{solution}

\begin{problem}
如图，在四面体 \(ABCD\) 中，平面 \(ABC \perp\) 平面 \(ACD\)，\(AB \perp BC\)，\(AC = AD = 2\)，\(BC = CD = 1\).
\begin{enumerate}
\item 求四面体 \(ABCD\) 的体积；
\item 求二面角 \(C - AB - D\) 的平面角的正切值.
\end{enumerate}
\bitmapfigure[max width=0.42\linewidth]{img/2011/chongqing_liberal/q20_fig1.png}
\end{problem}

\begin{solution}
由 \(AB\perp BC\)、\(AC=2\)、\(BC=1\)，得
\[
AB=\sqrt{AC^{2}-BC^{2}}=\sqrt3
\]

建立空间直角坐标系，使 \(B=(0,0,0)\)，\(A=(\sqrt3,0,0)\)，\(C=(0,1,0)\). 于是平面 \(ABC\) 为 \(z=0\)，且
\(S_{\triangle ABC}=\frac12\cdot\sqrt3\cdot1=\frac{\sqrt3}{2}\).

平面 \(ACD\) 垂直于 \(ABC\) 且交线为 \(AC\)，所以它是过 \(AC\) 的竖直平面，方程为
\(x+\sqrt3y=\sqrt3\). 设 \(D=(x,y,z)\)，由 \(AD=2\)、\(CD=1\) 得
\((x-\sqrt3)^{2}+y^{2}+z^{2}=4\)，\(x^{2}+(y-1)^{2}+z^{2}=1\).
两式相减得 \(y-\sqrt3x=\frac12\). 联立 \(x+\sqrt3y=\sqrt3\)，解得
\(x=\frac{\sqrt3}{8}\)，\(y=\frac78\). 再代入 \(CD^{2}=1\)，得
\(z^{2}=1-\frac{3}{64}-\frac{1}{64}=\frac{15}{16}\)，故点 \(D\) 到平面 \(ABC\) 的距离为 \(\frac{\sqrt{15}}4\).

\begin{enumerate}
\item 四面体体积为
\[
V=\frac13S_{\triangle ABC}\cdot\frac{\sqrt{15}}4
=\frac13\cdot\frac{\sqrt3}{2}\cdot\frac{\sqrt{15}}4
=\frac{\sqrt5}{8}
\]
\item 直线 \(AB\) 是 \(x\) 轴. 在平面 \(ABC\) 内，\(BC\) 垂直于 \(AB\)，可取其方向向量为 \((0,1,0)\)；在平面 \(ABD\) 内，取 \(BD\) 在垂直于 \(AB\) 方向上的分量 \((0,\frac78,\frac{\sqrt{15}}4)\). 这两个向量分别位于两个半平面内且都垂直于棱 \(AB\)，它们的夹角就是二面角 \(C-AB-D\) 的平面角. 故其正切值为
\[
\tan\theta=\frac{\frac{\sqrt{15}}4}{\frac78}=\frac{2\sqrt{15}}7
\]
\end{enumerate}
\end{solution}

\begin{problem}
如图，椭圆的中心为原点 \(O\)，离心率 \( e = \frac{\sqrt{2}}{2} \)，一条准线的方程是 \( x = 2\sqrt{2} \).
\begin{enumerate}
\item 求该椭圆的标准方程；
\item 设动点 \(P\) 满足： \(\overrightarrow{OP} = \overrightarrow{OM} + 2\overrightarrow{ON}\)，其中 \(M\)，\(N\) 是椭圆上的点，直线 \(OM\) 与 \(ON\) 的斜率之积为 \(-\frac{1}{2}\)，问：是否存在定点 \(F\)，使得 \(|PF|\) 与点 \(P\) 到直线 \(l: x = 2\sqrt{10}\) 的距离之比为定值？ 若存在，求 \(F\) 的坐标；若不存在，说明理由.
\end{enumerate}
\bitmapfigure[max width=0.42\linewidth]{img/2011/chongqing_liberal/q21_fig1.png}
\end{problem}

\begin{solution}
\begin{enumerate}
\item 由 \(e=\frac ca=\frac{\sqrt2}{2}\)，以及右准线 \(x=\frac ae=2\sqrt2\)，得
\(a=2\). 于是 \(c=ae=\sqrt2\)，再由 \(b^{2}=a^{2}-c^{2}\) 得 \(b^{2}=2\). 所以椭圆的标准方程为
\[
\frac{x^{2}}4+\frac{y^{2}}2=1
\]
\item 设 \(M=(x_{1},y_{1})\)，\(N=(x_{2},y_{2})\)，\(P=(x,y)\). 由
\(\overrightarrow{OP}=\overrightarrow{OM}+2\overrightarrow{ON}\)，得
\(x=x_{1}+2x_{2}\)，\(y=y_{1}+2y_{2}\). 又因为 \(M\)，\(N\) 在椭圆上，
\(x_{1}^{2}+2y_{1}^{2}=4\)，\(x_{2}^{2}+2y_{2}^{2}=4\).
直线 \(OM\)、\(ON\) 的斜率之积为 \(-\frac12\)，所以
\(\frac{y_{1}}{x_{1}}\cdot\frac{y_{2}}{x_{2}}=-\frac12\)，即
\(x_{1}x_{2}+2y_{1}y_{2}=0\). 因此
\[
x^{2}+2y^{2}=(x_{1}+2x_{2})^{2}+2(y_{1}+2y_{2})^{2}
\]
再利用上面的三个等式，
\[
x^{2}+2y^{2}=(x_{1}^{2}+2y_{1}^{2})+4(x_{2}^{2}+2y_{2}^{2})+4(x_{1}x_{2}+2y_{1}y_{2})=4+16+0=20
\]
故 \(P\) 在椭圆 \(\frac{x^{2}}{20}+\frac{y^{2}}{10}=1\) 上. 该椭圆的半长轴为 \(2\sqrt5\)，焦距参数为
\(C=\sqrt{20-10}=\sqrt{10}\)，右焦点为 \(F=(\sqrt{10},0)\)，离心率为
\(e_{1}=\frac{\sqrt{10}}{2\sqrt5}=\frac{\sqrt2}{2}\). 其右准线为
\(x=\frac{(2\sqrt5)^{2}}{\sqrt{10}}=2\sqrt{10}\)，恰为直线 \(l\). 由椭圆的焦点准线定义，对任意这样的 \(P\)，
\[
\frac{|PF|}{\operatorname{dist}(P,l)}=e_{1}=\frac{\sqrt2}{2}
\]
所以存在定点 \(F=(\sqrt{10},0)\)，定值为 \(\frac{\sqrt2}{2}\).
\end{enumerate}
\end{solution}
